\documentclass{article}
\usepackage{amsmath, amssymb, amsthm}
\usepackage{algorithm}
\usepackage{algorithmic}
\usepackage{booktabs}
\usepackage{hyperref}
\usepackage{graphicx}
\usepackage{xcolor}

\title{CoMI-IRL: Contrastive Multi-Intention Inverse Reinforcement Learning\\
\large Technical Methodology and Implementation Details}
\author{}
\date{}

\begin{document}
\maketitle

\tableofcontents
\newpage

%==============================================================================
\section{Problem Formulation}
%==============================================================================

We address the problem of \textbf{multi-intention inverse reinforcement learning}, where an agent must:
\begin{enumerate}
    \item Learn to cluster expert demonstrations into distinct behavioral modes (intentions)
    \item Train intention-specific reward functions and policies via IRL
    \item Adapt to novel behavioral modes encountered online without forgetting previously learned intentions
\end{enumerate}

\subsection{Formal Setting}

Let $\mathcal{D} = \{\tau_1, \tau_2, \ldots, \tau_N\}$ be a dataset of expert trajectories, where each trajectory $\tau_i = \{(s_t^i, a_t^i)\}_{t=1}^{T_i}$ is a sequence of state-action pairs. Each trajectory is generated by one of $K$ unknown behavioral modes (intentions), and our goal is to:

\begin{enumerate}
    \item \textbf{Cluster} trajectories by their underlying intention without access to ground-truth labels
    \item \textbf{Learn} intention-specific policies $\pi_k$ for each discovered cluster $k$
    \item \textbf{Detect} novel intentions when new trajectories arrive that don't match existing clusters
    \item \textbf{Adapt} by training new policies for novel intentions while preserving existing ones
\end{enumerate}

\subsection{Continual Learning Setup}

We formalize the continual setting as follows:
\begin{itemize}
    \item \textbf{Baseline Phase}: Observe trajectories from $K_{\text{seen}}$ modes: $\mathcal{D}_{\text{seen}} = \{\tau_i : \text{mode}(\tau_i) \in \mathcal{M}_{\text{seen}}\}$
    \item \textbf{Online Phase}: Receive new trajectories $\mathcal{D}_{\text{online}}$ that may contain:
    \begin{itemize}
        \item Trajectories from seen modes $\mathcal{M}_{\text{seen}}$
        \item Trajectories from novel modes $\mathcal{M}_{\text{novel}}$ where $\mathcal{M}_{\text{novel}} \cap \mathcal{M}_{\text{seen}} = \emptyset$
    \end{itemize}
\end{itemize}

The challenge is to correctly identify which online trajectories belong to existing clusters versus novel clusters, without access to ground-truth mode labels.

%==============================================================================
\section{Architecture: Behavior Encoder}
%==============================================================================

\subsection{Overview}

We employ a transformer-based \textbf{Behavior Encoder} that maps variable-length trajectories to fixed-dimensional embeddings. The encoder processes interleaved state-action sequences and produces a trajectory-level representation suitable for clustering.

\subsection{Input Representation}

For a trajectory $\tau = \{(s_t, a_t)\}_{t=1}^{T}$, we construct the input as follows:

\subsubsection{Random Fourier Features (RFF)}

To capture high-frequency patterns in continuous state and action spaces, we apply Random Fourier Features~\cite{rahimi2007random}:

\begin{equation}
    \phi_{\text{RFF}}(x) = \sqrt{\frac{2}{m}} \begin{bmatrix} \cos(Wx + b) \\ \sin(Wx + b) \end{bmatrix}
\end{equation}

where:
\begin{itemize}
    \item $W \in \mathbb{R}^{m \times d}$ with entries sampled from $\mathcal{N}(0, \sigma^2)$
    \item $b \in \mathbb{R}^{m}$ sampled uniformly from $[0, 2\pi]$
    \item $\sigma$ controls the bandwidth (we use $\sigma_{\text{state}} = 0.001$ for Walker2d/Pusher, $\sigma_{\text{state}} = 0.01$ for Reacher)
    \item $m_{\text{state}} = 64$, $m_{\text{action}} = 32$
\end{itemize}

\textbf{Justification}: RFF provides an explicit feature map approximation to the Gaussian RBF kernel, enabling the encoder to capture fine-grained behavioral differences that would be smoothed out by standard linear projections. This is particularly important for continuous control tasks where subtle differences in action timing distinguish behavioral modes.

\subsubsection{Temporal Convolution for Actions}

Actions are processed through a 1D temporal convolution before RFF:

\begin{equation}
    a'_t = \text{Conv1D}(a_{t-w:t+w})
\end{equation}

where $w$ is the kernel half-width. This captures local action dynamics and temporal correlations.

\textbf{Justification}: Actions in continuous control are temporally correlated (e.g., smooth motor commands). Temporal convolution captures these patterns, improving discrimination between modes with different action dynamics.

\subsection{Transformer Architecture}

The processed state-action tokens are fed to a transformer encoder:

\begin{equation}
    H = \text{TransformerEncoder}([\text{CLS}; \phi(s_1, a_1); \ldots; \phi(s_T, a_T)])
\end{equation}

Architecture details:
\begin{itemize}
    \item Embedding dimension: $d_{\text{emb}} = 32$
    \item Number of attention heads: $n_{\text{heads}} = 4$
    \item Number of layers: $n_{\text{layers}} = 2$
    \item Hidden dimension: $d_{\text{hid}} = 1024$
    \item Dropout: $p = 0.1$
\end{itemize}

The \texttt{[CLS]} token's final representation serves as the trajectory embedding:

\begin{equation}
    z_\tau = H_{\text{CLS}} \in \mathbb{R}^{d_{\text{emb}}}
\end{equation}

%==============================================================================
\section{Training Objectives}
%==============================================================================

The encoder is trained with a multi-objective loss combining contrastive learning, mutual information maximization, and temporal segmentation:

\begin{equation}
    \mathcal{L}_{\text{total}} = \alpha \mathcal{L}_{\text{contrastive}} + \beta \mathcal{L}_{\text{infomax}} + \gamma \mathcal{L}_{\text{segment}} + \delta \mathcal{L}_{\text{stability}}
\end{equation}

\subsection{Contrastive Loss ($\mathcal{L}_{\text{contrastive}}$)}

We use NT-Xent (Normalized Temperature-scaled Cross-Entropy) loss~\cite{chen2020simple} to encourage similar trajectories to have similar embeddings:

\begin{equation}
    \mathcal{L}_{\text{contrastive}} = -\frac{1}{N} \sum_{i=1}^{N} \log \frac{\exp(\text{sim}(z_i, z_i^+) / \tau)}{\sum_{j=1}^{N} \exp(\text{sim}(z_i, z_j) / \tau)}
\end{equation}

where:
\begin{itemize}
    \item $z_i^+$ is an augmented version of trajectory $i$ (temporal jittering, action noise)
    \item $\text{sim}(z_i, z_j) = \frac{z_i^\top z_j}{\|z_i\| \|z_j\|}$ is cosine similarity
    \item $\tau$ is the temperature parameter
\end{itemize}

\textbf{Justification}: Contrastive learning creates an embedding space where trajectories from the same behavioral mode cluster together, enabling unsupervised mode discovery.

\subsection{InfoMax Loss ($\mathcal{L}_{\text{infomax}}$)}

We maximize mutual information between the trajectory embedding and local temporal features, inspired by Deep InfoMax~\cite{hjelm2018learning}:

\begin{equation}
    \mathcal{L}_{\text{infomax}} = -\mathbb{E}\left[\log \sigma(f(z_\tau, h_t)) + \sum_{t' \neq t} \log(1 - \sigma(f(z_\tau, h_{t'}^-)))\right]
\end{equation}

where $f$ is a discriminator network and $h_t$ are local temporal representations.

\textbf{Justification}: InfoMax encourages the global trajectory embedding to preserve information about local temporal patterns, preventing the encoder from collapsing to trivial solutions.

\subsection{Segmentation Loss ($\mathcal{L}_{\text{segment}}$)}

For trajectories with known temporal structure, we encourage the encoder to capture segment boundaries:

\begin{equation}
    \mathcal{L}_{\text{segment}} = \text{BCE}(\hat{b}_t, b_t)
\end{equation}

where $b_t$ indicates segment boundaries and $\hat{b}_t$ are predicted boundary probabilities.

\subsection{Stability Loss ($\mathcal{L}_{\text{stability}}$) — Finetuning Only}

During finetuning, we add a stability term inspired by Elastic Weight Consolidation~\cite{kirkpatrick2017overcoming} to prevent catastrophic forgetting:

\begin{equation}
    \mathcal{L}_{\text{stability}} = \|z_\tau^{\text{new}} - z_\tau^{\text{old}}\|_2^2
\end{equation}

where $z_\tau^{\text{old}}$ is the embedding from the pre-finetuning encoder.

\textbf{Justification}: Without stability regularization, the encoder may catastrophically forget representations for seen modes when adapting to novel modes. This is especially problematic for continuous manifolds where modes overlap.

\subsection{Default Hyperparameters}

\begin{table}[h]
\centering
\begin{tabular}{lcc}
\toprule
\textbf{Parameter} & \textbf{Symbol} & \textbf{Default Value} \\
\midrule
Contrastive weight & $\alpha$ & 0.5 \\
InfoMax weight & $\beta$ & 1.0 \\
Segmentation weight & $\gamma$ & 0.5 \\
Stability weight & $\delta$ & 1.0 \\
Learning rate & $\eta$ & 0.0001 \\
Pretraining epochs & - & 50 \\
Finetuning epochs & - & 50 \\
\bottomrule
\end{tabular}
\caption{Training hyperparameters}
\end{table}

%==============================================================================
\section{Data Preprocessing}
%==============================================================================

\subsection{Normalization}

We apply quantile-normal scaling to both states and actions:

\begin{equation}
    x_{\text{scaled}} = \Phi^{-1}\left(F_{\text{empirical}}(x)\right)
\end{equation}

where $\Phi^{-1}$ is the inverse standard normal CDF and $F_{\text{empirical}}$ is the empirical CDF fitted on training data.

\textbf{Justification}: Quantile normalization is robust to outliers and produces approximately Gaussian marginals, which improves transformer training stability.

\textbf{Important}: Scalers are fitted \textbf{only on seen data} during baseline training. Online data is transformed using the same fitted scalers with clipping to handle out-of-distribution values.

%==============================================================================
\section{Phase 1: Baseline Clustering (Graph-Based Leiden)}
%==============================================================================

\subsection{Overview}

We employ graph-based community detection for clustering trajectory embeddings, using the Leiden algorithm~\cite{traag2019louvain} which provides guarantees on community connectivity that the earlier Louvain algorithm lacks.

\subsection{Graph Construction}

Given trajectory embeddings $\{z_i\}_{i=1}^{N}$, we construct a $k$-NN graph:

\begin{equation}
    \mathcal{G} = (V, E), \quad E = \{(i, j) : j \in \text{kNN}(i) \text{ or } i \in \text{kNN}(j)\}
\end{equation}

Edge weights are computed using cosine similarity:

\begin{equation}
    w_{ij} = \frac{z_i^\top z_j}{\|z_i\| \|z_j\|}
\end{equation}

We symmetrize the graph by including an edge $(i,j)$ if either $j \in \text{kNN}(i)$ or $i \in \text{kNN}(j)$.

\subsection{Automatic $k$ Selection via Joint Optimization}

We automatically select the optimal $k$ and resolution $\gamma$ by sweeping over a grid and evaluating clustering quality:

\begin{algorithm}[H]
\caption{Joint $k$-Resolution Selection}
\begin{algorithmic}[1]
\STATE \textbf{Input}: Embeddings $Z$, range $[k_{\min}, k_{\max}]$, resolution set $\Gamma$
\STATE \textbf{Output}: Optimal $(k^*, \gamma^*)$, cluster labels
\STATE
\STATE // Step 1: Check for isolated components
\FOR{$k \in \{15, 30, 50\}$}
    \STATE Build $k$-NN graph $\mathcal{G}_k$
    \STATE Compute connected components
    \IF{multiple significant components exist}
        \STATE Use connected components as clusters
        \RETURN labels, $k$
    \ENDIF
\ENDFOR
\STATE
\STATE // Step 2: Joint k-resolution sweep
\STATE $\mathcal{R} \leftarrow \emptyset$ \COMMENT{Results collection}
\FOR{$k \in \text{linspace}(k_{\min}, k_{\max}, 15)$}
    \STATE Build $k$-NN graph $\mathcal{G}_k$
    \FOR{$\gamma \in \Gamma$}
        \STATE Run Leiden clustering with resolution $\gamma$
        \STATE Compute silhouette score $s$, modularity $Q$
        \STATE $\mathcal{R} \leftarrow \mathcal{R} \cup \{(k, \gamma, s, Q, \text{labels})\}$
    \ENDFOR
\ENDFOR
\STATE
\STATE // Step 3: Selection via stability
\STATE Group results by $(n_{\text{clusters}})$
\STATE For each group, compute stability = count of configurations
\STATE Select $(k^*, \gamma^*)$ from most stable group with best silhouette
\RETURN labels, $(k^*, \gamma^*)$
\end{algorithmic}
\end{algorithm}

\subsection{Why Leiden over HDBSCAN?}

We chose Leiden community detection over HDBSCAN~\cite{campello2013density} for several reasons:

\begin{table}[h]
\centering
\begin{tabular}{p{6cm}p{6cm}}
\toprule
\textbf{HDBSCAN Limitation} & \textbf{Leiden Advantage} \\
\midrule
Struggles with continuous manifolds where modes blend & Graph structure captures local neighborhood relationships \\
Density-based, sensitive to varying cluster densities & Resolution parameter allows multi-scale community detection \\
No principled way to handle cosine similarity & Native support for weighted graphs with any similarity \\
Cluster assignments can be unstable & Guarantees well-connected communities \\
\bottomrule
\end{tabular}
\caption{Comparison of HDBSCAN vs Leiden for trajectory clustering}
\end{table}

\subsection{Leiden Community Detection}

The Leiden algorithm optimizes modularity with resolution parameter $\gamma$:

\begin{equation}
    Q = \frac{1}{2m} \sum_{ij} \left[ A_{ij} - \gamma \frac{k_i k_j}{2m} \right] \delta(c_i, c_j)
\end{equation}

where:
\begin{itemize}
    \item $A_{ij}$ = adjacency matrix (edge weights)
    \item $k_i$ = degree of node $i$
    \item $m$ = total edge weight
    \item $c_i$ = community assignment of node $i$
    \item $\gamma$ = resolution parameter (higher values yield more clusters)
\end{itemize}

The resolution parameter $\gamma$ controls granularity:
\begin{itemize}
    \item $\gamma < 1$: Fewer, larger communities
    \item $\gamma = 1$: Standard modularity optimization
    \item $\gamma > 1$: More, smaller communities
\end{itemize}

\subsection{Optional: Behavioral Feature Augmentation}

For environments with continuous manifolds (e.g., Walker2d), we optionally augment the graph with behavioral features derived from the policy Jacobian. This is based on the intuition that trajectories from the same behavioral mode should have similar state-action mappings.

\subsubsection{Jacobian Computation}

We estimate the local policy Jacobian via finite differences:

\begin{equation}
    J_\tau = \mathbb{E}_{(s,a) \sim \tau}\left[\frac{\partial a}{\partial s}\right] \approx \frac{1}{T-1} \sum_{t=1}^{T-1} \frac{a_{t+1} - a_t}{s_{t+1} - s_t + \epsilon}
\end{equation}

The Jacobian matrix is flattened to a feature vector for each trajectory.

\subsubsection{Edge Reweighting}

We incorporate Jacobian similarity into edge weights:

\begin{equation}
    w_{ij}^{\text{aug}} = (1 - \alpha_{\text{behav}}) \cdot w_{ij}^{\text{embed}} + \alpha_{\text{behav}} \cdot w_{ij}^{\text{jacobian}}
\end{equation}

where $\alpha_{\text{behav}} \in [0, 1]$ controls the behavioral feature contribution.

\subsubsection{Redundancy Test}

Before using behavioral features, we compute the Pearson and Spearman correlation between embedding-based and Jacobian-based edge weights:

\begin{equation}
    \rho = \text{corr}(w^{\text{embed}}, w^{\text{jacobian}})
\end{equation}

If $|\rho| > 0.7$, behavioral features are redundant (already captured by the encoder) and we skip them. This prevents adding noise without benefit.

\subsection{Cluster Registry}

After clustering, we build a registry storing cluster metadata:

\begin{equation}
    \text{Registry}[k] = \left\{
    \begin{array}{l}
        \mu_k = \frac{1}{|C_k|} \sum_{i \in C_k} z_i \quad \text{(centroid)} \\
        r_k^{95} = \text{quantile}_{95}\left(\{d_{\cos}(z_i, \mu_k)\}_{i \in C_k}\right) \quad \text{(radius)} \\
        n_k = |C_k| \quad \text{(count)}
    \end{array}
    \right.
\end{equation}

where $d_{\cos}(z_i, \mu_k) = 1 - \frac{z_i^\top \mu_k}{\|z_i\| \|\mu_k\|}$ is the cosine distance.

%==============================================================================
\section{Phase 1.5: Baseline IRL Training}
%==============================================================================

For each discovered cluster $k$, we train an intention-specific policy using Inverse Reinforcement Learning.

\subsection{Supported IRL Algorithms}

\begin{itemize}
    \item \textbf{GAIL} (Generative Adversarial Imitation Learning)~\cite{ho2016generative}: Learns a discriminator to distinguish expert from learner state-action pairs
    \item \textbf{AIRL} (Adversarial IRL)~\cite{fu2018learning}: Recovers a disentangled reward function robust to dynamics changes
    \item \textbf{SQIL} (Soft Q Imitation Learning)~\cite{reddy2020sqil}: Assigns reward +1 to expert transitions and 0 to learner transitions
\end{itemize}

\subsection{Training Procedure}

For each cluster $k$:
\begin{enumerate}
    \item Extract trajectories: $\mathcal{D}_k = \{\tau_i : \ell_i = k\}$
    \item Train IRL agent for $T_{\text{RL}}$ timesteps (250K for Reacher/Pusher, 1M for Walker2d)
    \item Evaluate learned policy against expert demonstrations
\end{enumerate}

\subsection{Parallel Training}

We support parallel IRL training across clusters using CPU multiprocessing:

\begin{equation}
    \text{Workers} = \min(\text{num\_clusters}, \text{CPU\_cores} - 1)
\end{equation}

%==============================================================================
\section{Phase 2: Encoder Finetuning}
%==============================================================================

When online trajectories $\mathcal{D}_{\text{online}}$ arrive, we finetune the encoder on the combined dataset:

\begin{equation}
    \mathcal{D}_{\text{combined}} = \mathcal{D}_{\text{seen}} \cup \mathcal{D}_{\text{online}}
\end{equation}

\subsection{Finetuning Loss}

The finetuning loss includes a stability term to prevent catastrophic forgetting:

\begin{equation}
    \mathcal{L}_{\text{finetune}} = \alpha \mathcal{L}_{\text{contrastive}} + \beta \mathcal{L}_{\text{infomax}} + \gamma \mathcal{L}_{\text{segment}} + \delta \mathcal{L}_{\text{stability}}
\end{equation}

The stability loss penalizes drift in embeddings for seen trajectories:

\begin{equation}
    \mathcal{L}_{\text{stability}} = \frac{1}{|\mathcal{D}_{\text{seen}}|} \sum_{\tau \in \mathcal{D}_{\text{seen}}} \|f_\theta(\tau) - f_{\theta_0}(\tau)\|_2^2
\end{equation}

where $\theta_0$ are the pre-finetuning parameters.

\textbf{Justification}: Without stability regularization, the encoder may catastrophically forget representations for seen modes when adapting to novel modes. This is especially problematic for continuous manifolds where modes overlap.

%==============================================================================
\section{Phase 3: Two-Stage Target-Aware Clustering}
%==============================================================================

This is the \textbf{key methodological contribution} for handling continuous behavioral manifolds. After finetuning, the embedding space may have shifted significantly, causing baseline cluster labels to no longer correspond to the correct modes.

\subsection{The Problem: Embedding Drift}

After finetuning, we observe:
\begin{itemize}
    \item \textbf{High drift}: Cluster centroids move significantly in embedding space
    \item \textbf{Label-mode mismatch}: Trajectories with baseline label ``Cluster 0'' may now be closer to what was ``Cluster 1''
    \item \textbf{Mixing}: Seen modes may partially overlap with novel modes
\end{itemize}

A naive approach of anchoring to baseline cluster labels fails because those labels were computed in a different (pre-finetuning) embedding space.

\subsection{Solution: Two-Stage Target-Aware Clustering}

We propose a two-stage approach that leverages knowledge of the expected number of baseline clusters:

\begin{algorithm}[H]
\caption{Two-Stage Target-Aware Clustering}
\begin{algorithmic}[1]
\STATE \textbf{Input}: Finetuned embeddings $Z_{\text{all}}$, number of seen trajectories $N_{\text{seen}}$, expected baseline clusters $K_{\text{baseline}}$
\STATE \textbf{Output}: Cluster labels, novel cluster IDs

\STATE
\STATE \textbf{// ========== STAGE 1: Target-Aware Re-clustering ==========}
\STATE $Z_{\text{seen}} \leftarrow Z_{\text{all}}[1:N_{\text{seen}}]$
\STATE
\STATE // Step 1a: Check for isolated components
\FOR{$k \in \{15, 30, 50\}$}
    \STATE Build $k$-NN graph on $Z_{\text{seen}}$
    \STATE Count significant connected components
    \IF{num\_components $= K_{\text{baseline}}$}
        \STATE Use components as recovered clusters
        \STATE \textbf{goto} Stage 2
    \ENDIF
\ENDFOR
\STATE
\STATE // Step 1b: Target-aware k-resolution sweep
\STATE $\mathcal{R} \leftarrow \emptyset$, $\mathcal{E} \leftarrow \emptyset$ \COMMENT{All results, exact matches}
\FOR{$k \in \text{linspace}(k_{\min}, k_{\max}, 15)$}
    \FOR{$\gamma \in \{0.01, 0.025, 0.05, \ldots, 2.0\}$}
        \STATE Run Leiden clustering
        \STATE $n_c \leftarrow$ number of valid clusters
        \STATE Compute silhouette $s$
        \STATE $\mathcal{R} \leftarrow \mathcal{R} \cup \{(k, \gamma, n_c, s)\}$
        \IF{$n_c = K_{\text{baseline}}$}
            \STATE $\mathcal{E} \leftarrow \mathcal{E} \cup \{|\mathcal{R}|\}$ \COMMENT{Track exact matches}
        \ENDIF
    \ENDFOR
\ENDFOR
\STATE
\STATE // Step 1c: Selection strategy
\IF{$\mathcal{E} \neq \emptyset$}
    \STATE Select from $\mathcal{E}$ with best silhouette \COMMENT{Exact match}
\ELSE
    \STATE Score $= -2 \cdot |n_c - K_{\text{baseline}}| + s$
    \STATE Select configuration with best score \COMMENT{Closest to target}
\ENDIF
\STATE
\STATE Compute recovered centroids $\{\mu_k^{\text{rec}}\}$ and radii $\{r_k^{\text{rec}}\}$

\STATE
\STATE \textbf{// ========== STAGE 2: Novel Detection via Anchoring ==========}
\STATE $Z_{\text{online}} \leftarrow Z_{\text{all}}[N_{\text{seen}}+1:]$
\FOR{each online trajectory $z_i$}
    \STATE $d_k \leftarrow d_{\cos}(z_i, \mu_k^{\text{rec}})$ for all recovered clusters $k$
    \IF{$\min_k d_k \leq r_k^{\text{rec}} \cdot \theta_{\text{novelty}}$}
        \STATE Assign to nearest cluster: $\ell_i \leftarrow \arg\min_k d_k$
    \ELSE
        \STATE Mark as novel candidate
    \ENDIF
\ENDFOR

\STATE
\STATE \textbf{// ========== STAGE 2b: Sub-cluster Novel Candidates ==========}
\IF{novel candidates $\geq$ min\_cluster\_size}
    \STATE Build $k$-NN graph on novel candidates
    \STATE Compute connected components
    \IF{multiple components}
        \STATE Each component $\rightarrow$ separate novel cluster
    \ELSE
        \STATE Run Leiden to find sub-clusters
    \ENDIF
\ENDIF

\STATE
\RETURN labels, novel\_cluster\_ids
\end{algorithmic}
\end{algorithm}

\subsection{Stage 1: Target-Aware Recovery of Baseline Structure}

\textbf{Key Insight}: We know the expected number of clusters $K_{\text{baseline}}$ from Phase 1. Even though the embedding space has shifted, the relative structure of seen modes is often preserved. Mode 0 and Mode 1 remain separable; they've just moved in the embedding space.

\subsubsection{Target-Aware Selection Strategy}

Given the target cluster count $K_{\text{baseline}}$, we prioritize configurations that match this target:

\begin{enumerate}
    \item \textbf{Exact Matches}: If any $(k, \gamma)$ configuration produces exactly $K_{\text{baseline}}$ clusters, we select among these using silhouette score as tiebreaker
    \item \textbf{Closest Match}: If no exact match exists, we score configurations as:
    \begin{equation}
        \text{score}(k, \gamma) = -2 \cdot |n_{\text{clusters}} - K_{\text{baseline}}| + s_{\text{silhouette}}
    \end{equation}
    This heavily penalizes deviation from target while still rewarding cluster quality.
\end{enumerate}

\subsubsection{Output of Stage 1}

Stage 1 produces:
\begin{itemize}
    \item \textbf{Recovered centroids} $\mu_k^{\text{rec}}$: New cluster centers in the finetuned embedding space
    \item \textbf{Recovered radii} $r_k^{\text{rec}}$: 95th percentile of within-cluster cosine distances, expanded by factor 1.2
\end{itemize}

\subsection{Stage 2: Novel Detection via Distance Thresholding}

For each online trajectory, we compute its distance to all recovered centroids:

\begin{equation}
    d_k(z) = 1 - \frac{z^\top \mu_k^{\text{rec}}}{\|z\| \|\mu_k^{\text{rec}}\|}
\end{equation}

A trajectory is classified as:
\begin{itemize}
    \item \textbf{Existing mode}: if $\exists k : d_k(z) \leq \theta_{\text{novelty}} \cdot r_k^{\text{rec}}$
    \item \textbf{Novel candidate}: otherwise
\end{itemize}

The novelty threshold $\theta_{\text{novelty}}$ controls sensitivity:
\begin{itemize}
    \item Lower values (e.g., 0.05): Strict — more trajectories marked as novel
    \item Higher values (e.g., 1.5): Lenient — only clear outliers marked as novel
\end{itemize}

We use $\theta_{\text{novelty}} = 0.05$ for Walker2d-v4 (continuous manifold) and $\theta_{\text{novelty}} = 0.1$ for Reacher/Pusher (isolated clusters).

\subsection{Stage 2b: Novel Candidate Sub-Clustering}

Novel candidates are grouped into clusters via:
\begin{enumerate}
    \item \textbf{Connected component analysis}: If novel candidates form multiple disconnected components in the $k$-NN graph, each component becomes a separate novel cluster
    \item \textbf{Leiden clustering}: If a single connected component, apply Leiden to potentially discover sub-clusters
    \item \textbf{Single cluster}: If too few candidates or single cohesive group, assign all to one novel cluster
\end{enumerate}

\subsection{Why Two-Stage Target-Aware Clustering Works}

\begin{table}[h]
\centering
\begin{tabular}{p{5.5cm}p{6.5cm}}
\toprule
\textbf{Problem with Naive Approaches} & \textbf{Two-Stage Target-Aware Solution} \\
\midrule
Anchoring to stale baseline labels that no longer match modes after embedding drift & Re-clusters seen data in new embedding space using target-aware selection \\
No guidance on expected cluster count leads to over/under-segmentation & Uses known $K_{\text{baseline}}$ to guide k-resolution selection \\
Cannot distinguish mode boundary cases from true novelty & Uses recovered radii as calibrated novelty thresholds \\
Joint clustering of seen+online may merge novel with existing & Stage 2 anchors online data to recovered (not stale) clusters \\
\bottomrule
\end{tabular}
\caption{Problems solved by two-stage target-aware clustering}
\end{table}

%==============================================================================
\section{Phase 4: Novel Cluster Analysis and IRL Training}
%==============================================================================

\subsection{Novel Cluster Validation}

For each detected novel cluster, we analyze:

\begin{enumerate}
    \item \textbf{Size}: Number of trajectories assigned
    \item \textbf{Purity}: Fraction of trajectories from the dominant true mode (for validation with ground truth)
    \item \textbf{Mode distribution}: Histogram of true modes within the cluster
    \item \textbf{Expected novel check}: Whether the dominant mode is in $\mathcal{M}_{\text{unseen}}$
\end{enumerate}

\subsection{IRL Training for Novel Clusters}

Novel clusters undergo the same IRL training procedure as baseline clusters:

\begin{enumerate}
    \item Extract trajectories from novel cluster
    \item Train intention-specific policy (GAIL/AIRL/SQIL)
    \item Evaluate against expert demonstrations
\end{enumerate}

%==============================================================================
\section{Evaluation Metrics}
%==============================================================================

\subsection{Clustering Quality Metrics}

\subsubsection{Normalized Mutual Information (NMI)}

\begin{equation}
    \text{NMI}(Y, \hat{Y}) = \frac{2 \cdot I(Y; \hat{Y})}{H(Y) + H(\hat{Y})}
\end{equation}

where $I$ is mutual information and $H$ is entropy. Range: $[0, 1]$, higher is better.

\subsubsection{Adjusted Rand Index (ARI)}

\begin{equation}
    \text{ARI} = \frac{\text{RI} - \mathbb{E}[\text{RI}]}{\max(\text{RI}) - \mathbb{E}[\text{RI}]}
\end{equation}

Measures agreement between true and predicted clusterings, adjusted for chance. Range: $[-1, 1]$, higher is better.

\subsubsection{Silhouette Score}

\begin{equation}
    s(i) = \frac{b(i) - a(i)}{\max(a(i), b(i))}
\end{equation}

where $a(i)$ is mean intra-cluster distance and $b(i)$ is mean nearest-cluster distance. Range: $[-1, 1]$, higher is better.

\subsection{Novel Detection Metrics}

These are the \textbf{primary metrics} for evaluating the continual learning capability:

\subsubsection{Precision}

\begin{equation}
    \text{Precision} = \frac{\text{TP}}{\text{TP} + \text{FP}}
\end{equation}

Of trajectories predicted as novel, what fraction are truly from unseen modes?

\subsubsection{Recall}

\begin{equation}
    \text{Recall} = \frac{\text{TP}}{\text{TP} + \text{FN}}
\end{equation}

Of truly unseen trajectories, what fraction were correctly identified as novel?

\subsubsection{F1 Score}

\begin{equation}
    \text{F1} = \frac{2 \cdot \text{Precision} \cdot \text{Recall}}{\text{Precision} + \text{Recall}}
\end{equation}

Harmonic mean of precision and recall.

\subsubsection{Confusion Matrix Elements}

\begin{itemize}
    \item \textbf{TP} (True Positives): Unseen trajectories correctly placed in novel clusters
    \item \textbf{FP} (False Positives): Seen trajectories incorrectly placed in novel clusters
    \item \textbf{FN} (False Negatives): Unseen trajectories incorrectly assigned to baseline clusters
    \item \textbf{TN} (True Negatives): Seen trajectories correctly kept in baseline clusters
\end{itemize}

\subsection{IRL Performance Metrics}

\begin{itemize}
    \item \textbf{Expert Reward Mean/Std}: Average episodic reward of expert demonstrations
    \item \textbf{Learner Reward Mean/Std}: Average episodic reward of learned policy
    \item \textbf{Performance Ratio}: $\frac{\text{Learner Reward}}{\text{Expert Reward}}$
\end{itemize}

%==============================================================================
\section{Environment-Specific Considerations}
%==============================================================================

\subsection{Walker2d-v4: Continuous Behavioral Manifold}

Walker2d presents a unique challenge: behavioral modes form a \textbf{continuous manifold} rather than isolated clusters.

\begin{itemize}
    \item \textbf{Mode 0}: Forward walking gait
    \item \textbf{Mode 1}: Backward walking gait  
    \item \textbf{Mode 2}: Intermediate/transition gait (acts as ``bridge'' between modes 0 and 1)
\end{itemize}

\textbf{Implications}:
\begin{itemize}
    \item Standard density-based clustering (HDBSCAN) fails to separate modes
    \item Behavioral features (Jacobian) provide marginal benefit due to gait similarity
    \item Two-stage target-aware clustering is \textbf{essential} to handle the continuous manifold
    \item Novel detection is more challenging at mode boundaries
    \item Lower novelty threshold ($\theta = 0.05$) prevents over-assignment to existing clusters
\end{itemize}

\subsection{Reacher-v4 and Pusher-v4: Isolated Clusters}

These environments have modes corresponding to different target locations, producing well-separated clusters:

\begin{itemize}
    \item 6 distinct modes (target positions arranged in a circle)
    \item Modes are geometrically isolated in embedding space
    \item Standard Leiden clustering works well without target-aware selection
    \item Alternating mode selection (even indices seen, odd indices unseen) ensures diversity
    \item Higher novelty threshold ($\theta = 0.1$) is sufficient
\end{itemize}

%==============================================================================
\section{Complete Pipeline Summary}
%==============================================================================

\begin{algorithm}[H]
\caption{CoMI-IRL: Complete Pipeline}
\begin{algorithmic}[1]
\STATE \textbf{Input}: Expert trajectories $\mathcal{D}$, seen modes $\mathcal{M}_{\text{seen}}$, unseen modes $\mathcal{M}_{\text{unseen}}$

\STATE
\STATE \textbf{// ===== Phase 0: Data Preparation =====}
\STATE Split $\mathcal{D}$ into $\mathcal{D}_{\text{seen}}$ and $\mathcal{D}_{\text{online}}$ by mode
\STATE Fit quantile-normal scalers on $\mathcal{D}_{\text{seen}}$
\STATE Apply scalers to all data (with clipping for online)

\STATE
\STATE \textbf{// ===== Phase 1: Baseline Encoder Training =====}
\STATE Train encoder $f_\theta$ on $\mathcal{D}_{\text{seen}}$ with $\mathcal{L} = \alpha\mathcal{L}_c + \beta\mathcal{L}_i + \gamma\mathcal{L}_s$
\STATE Extract embeddings $Z_{\text{seen}} = \{f_\theta(\tau) : \tau \in \mathcal{D}_{\text{seen}}\}$

\STATE
\STATE \textbf{// ===== Phase 1: Baseline Clustering (Leiden) =====}
\STATE Compute Jacobian redundancy test
\STATE Run joint k-resolution sweep with optional behavioral features
\STATE Obtain baseline labels $\ell_{\text{baseline}}$, registry $\mathcal{R}$, $K_{\text{baseline}}$ clusters

\STATE
\STATE \textbf{// ===== Phase 1.5: Baseline IRL =====}
\FOR{each cluster $k \in \mathcal{R}$}
    \STATE Train IRL agent $\pi_k$ on cluster trajectories
\ENDFOR

\STATE
\STATE \textbf{// ===== Phase 2: Encoder Finetuning =====}
\STATE Finetune $f_\theta$ on $\mathcal{D}_{\text{seen}} \cup \mathcal{D}_{\text{online}}$ with stability loss
\STATE Extract finetuned embeddings $Z_{\text{all}}$

\STATE
\STATE \textbf{// ===== Phase 3: Two-Stage Target-Aware Clustering =====}
\STATE \textbf{Stage 1}: Re-cluster $Z_{\text{seen}}$ with target $K_{\text{baseline}}$
\STATE \textbf{Stage 2}: Assign $Z_{\text{online}}$ using recovered anchors
\STATE Identify novel clusters $\mathcal{N}$

\STATE
\STATE \textbf{// ===== Phase 4: Novel Cluster Analysis \& IRL =====}
\STATE Validate novel clusters (purity, mode distribution)
\STATE Compute novel detection metrics (precision, recall, F1)
\FOR{each novel cluster $k \in \mathcal{N}$}
    \STATE Train IRL agent $\pi_k$ on novel cluster trajectories
\ENDFOR

\STATE
\RETURN All trained policies $\{\pi_k\}$, clustering metrics, novel detection metrics
\end{algorithmic}
\end{algorithm}

%==============================================================================
\section{Implementation Details}
%==============================================================================

\subsection{Software Dependencies}

\begin{itemize}
    \item \textbf{PyTorch}: Neural network training (encoder)
    \item \textbf{Stable-Baselines3}: RL/IRL algorithms (PPO, GAIL, AIRL)
    \item \textbf{imitation}: IRL algorithm implementations
    \item \textbf{scikit-learn}: Clustering metrics, preprocessing
    \item \textbf{leidenalg}: Leiden community detection (via igraph)
    \item \textbf{UMAP}: Dimensionality reduction for visualization
    \item \textbf{Gymnasium}: Environment interfaces
\end{itemize}

\subsection{Hardware Requirements}

\begin{itemize}
    \item GPU (CUDA) or Apple Silicon (MPS) for encoder training
    \item CPU parallelization for IRL training (one worker per cluster)
    \item Typical runtime: 2-4 hours for complete pipeline (depending on number of clusters and IRL timesteps)
\end{itemize}

\subsection{Reproducibility}

All random seeds are fixed for reproducibility:
\begin{itemize}
    \item PyTorch: \texttt{torch.manual\_seed(SEED)}
    \item NumPy: \texttt{np.random.seed(SEED)}
    \item Python: \texttt{random.seed(SEED)}
    \item UMAP: \texttt{random\_state=SEED}
    \item Leiden: \texttt{seed=SEED}
\end{itemize}

%==============================================================================
\section{Results Summary}
%==============================================================================

\subsection{Output CSV Format}

The pipeline outputs a CSV file containing:

\begin{table}[h]
\centering
\small
\begin{tabular}{ll}
\toprule
\textbf{Column} & \textbf{Description} \\
\midrule
seed & Random seed used \\
env\_name & Environment (Reacher-v4/Pusher-v4/Walker2d-v4) \\
cluster\_id & Cluster identifier (baseline ID or novel\_X) \\
stage & Training phase (baseline/finetuned\_novel) \\
NMI\_final & Normalized Mutual Information \\
ARI\_final & Adjusted Rand Index \\
Silhouette\_final & Silhouette Score \\
novel\_precision & Precision of novel detection \\
novel\_recall & Recall of novel detection \\
novel\_f1 & F1 score of novel detection \\
novel\_TP/FP/FN/TN & Confusion matrix elements \\
n\_novel\_clusters & Number of novel clusters detected \\
Expert Reward Mean & Mean expert episodic reward \\
Learner Reward Mean & Mean learned policy reward \\
\bottomrule
\end{tabular}
\caption{Output CSV columns}
\end{table}

\subsection{Typical Results (Walker2d-v4)}

\begin{table}[h]
\centering
\begin{tabular}{lccc}
\toprule
\textbf{Metric} & \textbf{Baseline (2 modes)} & \textbf{After Finetuning (3 modes)} \\
\midrule
NMI & 1.00 & 0.91 \\
ARI & 1.00 & 0.93 \\
Novel Precision & - & 1.00 \\
Novel Recall & - & 0.93 \\
Novel F1 & - & 0.96 \\
\bottomrule
\end{tabular}
\caption{Clustering and novel detection performance on Walker2d-v4}
\end{table}

%==============================================================================
\section{Related Work}
%==============================================================================

\subsection{Trajectory Clustering and Representation Learning}

Our encoder architecture builds on advances in sequence modeling and representation learning. The use of transformers for trajectory encoding follows the success of attention mechanisms in sequential decision-making tasks. Random Fourier Features~\cite{rahimi2007random} provide an efficient approximation to kernel methods, enabling capture of high-frequency behavioral patterns.

\subsection{Community Detection Algorithms}

We employ the Leiden algorithm~\cite{traag2019louvain}, which improves upon the widely-used Louvain algorithm by guaranteeing that all communities are well-connected. This property is essential for our application, where we need stable, coherent behavioral clusters. Unlike density-based methods such as HDBSCAN~\cite{campello2013density}, graph-based community detection naturally handles the cosine similarity metric and varying cluster densities present in trajectory embeddings.

\subsection{Inverse Reinforcement Learning}

Our IRL training leverages established algorithms including GAIL~\cite{ho2016generative}, which frames imitation as a game between a policy and a discriminator, and AIRL~\cite{fu2018learning}, which recovers disentangled reward functions. SQIL~\cite{reddy2020sqil} provides a simpler alternative that assigns binary rewards.

\subsection{Continual Learning}

The stability loss in our finetuning phase is inspired by Elastic Weight Consolidation~\cite{kirkpatrick2017overcoming}, which prevents catastrophic forgetting by regularizing important weights. Our approach differs by operating in embedding space rather than weight space, preserving the structure of learned representations.

\subsection{Contrastive Learning}

Our contrastive loss follows the NT-Xent formulation from SimCLR~\cite{chen2020simple}, adapted for trajectory data with appropriate augmentations. The InfoMax component draws from Deep InfoMax~\cite{hjelm2018learning}, encouraging global representations to capture local structure.

%==============================================================================
\begin{thebibliography}{99}

\bibitem{rahimi2007random}
A. Rahimi and B. Recht, ``Random Features for Large-Scale Kernel Machines,'' in \textit{Advances in Neural Information Processing Systems}, 2007.

\bibitem{traag2019louvain}
V.A. Traag, L. Waltman, and N.J. van Eck, ``From Louvain to Leiden: guaranteeing well-connected communities,'' \textit{Scientific Reports}, vol. 9, no. 1, pp. 1-12, 2019.

\bibitem{campello2013density}
R.J.G.B. Campello, D. Moulavi, and J. Sander, ``Density-Based Clustering Based on Hierarchical Density Estimates,'' in \textit{Pacific-Asia Conference on Knowledge Discovery and Data Mining}, pp. 160-172, 2013.

\bibitem{ho2016generative}
J. Ho and S. Ermon, ``Generative Adversarial Imitation Learning,'' in \textit{Advances in Neural Information Processing Systems}, pp. 4565-4573, 2016.

\bibitem{fu2018learning}
J. Fu, K. Luo, and S. Levine, ``Learning Robust Rewards with Adversarial Inverse Reinforcement Learning,'' in \textit{International Conference on Learning Representations}, 2018.

\bibitem{reddy2020sqil}
S. Reddy, A.D. Dragan, and S. Levine, ``SQIL: Imitation Learning via Reinforcement Learning with Sparse Rewards,'' in \textit{International Conference on Learning Representations}, 2020.

\bibitem{kirkpatrick2017overcoming}
J. Kirkpatrick, R. Pascanu, N. Rabinowitz, et al., ``Overcoming catastrophic forgetting in neural networks,'' \textit{Proceedings of the National Academy of Sciences}, vol. 114, no. 13, pp. 3521-3526, 2017.

\bibitem{chen2020simple}
T. Chen, S. Kornblith, M. Norouzi, and G. Hinton, ``A Simple Framework for Contrastive Learning of Visual Representations,'' in \textit{International Conference on Machine Learning}, pp. 1597-1607, 2020.

\bibitem{hjelm2018learning}
R.D. Hjelm, A. Fedorov, S. Lavoie-Marchildon, et al., ``Learning deep representations by mutual information estimation and maximization,'' in \textit{International Conference on Learning Representations}, 2019.

\end{thebibliography}

\end{document}